\cvsection{Progetti Personali}
\smallskip
\cvproject{Kakebo}{App Android · 2026}{\cvtag{Flutter}\cvtag{Google ML Kit}\cvtag{OCR}\cvtag{Information Extraction}\cvtag{On-device AI}}{\href{https://github.com/AlessioBarbanti/Kakebo}{Kakebo}}
{\small Lettura degli scontrini \textbf{interamente sul telefono}: OCR con Google ML Kit ed \textbf{estrazione di totale, data e negozio} dal testo riconosciuto, con le righe ricostruite dalla posizione nella foto e la correzione degli errori tipici dell'OCR; nessun dato lascia il dispositivo. L'app segue il \textit{kakebo}, il metodo giapponese del 1904 per annotare ogni spesa e riflettere a fine mese su come spendere meglio.}

\medskip
\cvproject{Pidgeon Survivor}{Gioco Android · 2026}{\cvtag{AI Agents}\cvtag{Prompt Engineering}\cvtag{Automated Testing}\cvtag{Image Generation}\cvtag{Godot}}{\href{https://github.com/AlessioBarbanti/PidgeonSurvivor}{PidgeonSurvivor}}
{\small Gioco 2D sviluppato con un \textbf{sistema di agenti AI} progettato su misura: sei subagenti in Claude~Code (bilanciamento, QA, revisione UX, direzione artistica, asset, analisi del codice) e flussi automatizzati che portano ogni modifica dall'apertura alla verifica su un telefono reale. Oltre 170 file di test automatici; asset grafici generati con Codex.}
